
Model repositories such as HuggingFace host collections of trained neural networks spanning diverse architectures. Predicting model quality from weights alone, without task-specific evaluation, remains challenging when model zoos contain heterogeneous architecture families (CNNs, Transformers, MLPs). Prior weight-space learning methods demonstrate success on homogeneous model collections---Neural Functional Transformers on single-family MNIST MLPs \citep{zhou2023meta}, or architecture-specific methods like DWSNets for convolutional networks---but do not address heterogeneous settings where multiple architecture types coexist.

We investigate whether architecture-specific tokenization followed by shared attention-based processing can enable weight-space learning across heterogeneous model families. The approach projects diverse weight types (convolutional kernels, Q/K/V attention matrices, fully-connected weight matrices) into a shared D=128 dimensional token space using family-specific linear projections, then applies permutation-equivariant Neural Functional Transformer (NFT) attention to learn quality-predictive representations.

We evaluate this compositional design through five validation components at 150-model proof-of-concept scale: (1) per-family ablation to test whether tokenization preserves architecture-specific signals, (2) clustering analysis to verify that shared processing maintains family structure, (3-4) comparisons against flat-weight and engineered-feature baselines, and (5) robustness validation across tokenization variants. The mechanism validation components pass their prespecified thresholds (5/5 components), though overall performance on the heterogeneous test set ($\rho$=0.294, n=30 test samples) indicates that larger-scale training may be required to reach target performance levels.

The paper is structured as follows: Section 2 reviews related work, Section 3 describes the method, Section 4 details the experimental setup, Section 5 presents results, Section 6 discusses findings and limitations, and Section 7 concludes.
